\documentclass[../../main.tex]{subfiles}

\begin{document}

\subsection{I neutrini nel Modello Standard}
  \label{subsec:il-neutrino-nel-modello-standard}
  Il Modello Standard è una teoria che classifica tutte le particelle
  fondamentali e descrive le leggi che regolano le loro interazioni.
  Secondo questa teoria i neutrini appartengono alla categoria dei
  leptoni: particelle con spin $1/2$, non soggette all'interazione forte,
  organizzate in 3 generazioni.\\
  \begin{equation*}
    \begin{pmatrix}
      \Pnue \\ \Pelectron
    \end{pmatrix}
    \begin{pmatrix}
      \Pnum \\ \Pmuon
    \end{pmatrix}
    \begin{pmatrix}
      \Pnut \\ \Ptauon
    \end{pmatrix}
  \end{equation*}\\
  I neutrini inoltre, al contrario degli altri leptoni, sono caratterizzati
  da massa e carica elettrica nulle, possono quindi interagire solo tramite
  interazione debole.\\
  Le interazioni deboli dei neutrini possono essere classificate in base alla
  natura delle particelle con le quali interagiscono, di seguito sono esposte le
  interazioni con i costituenti della materia.

  \paragraph{Scattering neutrino elettrone}
    Lo scattering neutrino elettrone riguarda solo leptoni liberi, è quindi
    possibile calcolarne esplicitamente l'ampiezza tramite le regole di Feynman.
    Queste interazioni possono cambiare la natura delle particelle coinvolte
    \begin{equation*}
      \Pnulepton + \Pelectron \rightarrow \Pnue + \Pleptonminus
    \end{equation*}
    e in tal caso sono dette quasi elastiche, altrimenti si definiscono
    elastiche e assumo la forma
    \begin{align*}
      \Pnulepton + \Pelectron \rightarrow \Pnulepton + \Pelectron \\
      \APnulepton + \Pelectron \rightarrow \APnulepton + \Pelectron
    \end{align*}
    Gli scattering elastici, lasciando invariate le particelle interagenti, non
    richiedono un'energia minima per avvenire, al contrario gli scattering quasi
    elastici richiedono che il neutrino abbia un'energia sufficiente (nel
    sistema di riferimento dell'elettrone) per creare il leptone carico
    corrispondente.

  \paragraph{Scattering neutrino nucleone}
    La struttura a quark dei nucleoni rende impossibile calcolare la sezione
    d'urto di queste interazioni tramite le regole di Feynman e richiede
    tecniche più avanzate, non trattate in questo elaborato.
    Queste interazioni si dividono in tre gruppi: elastici, quasi elastici e
    anelastici.
    I primi due casi sono analoghi agli omonimi già visti per gli scattering
    neutrino elettrone; si ricorda però che al fine di conservare la carica
    elettrica i neutrini possono interagire quasi elasticamente solo con
    neutroni mentre gli antineutrini solo con protoni
    \begin{align*}
      \Pnulepton + \Pneutron \rightarrow \Pleptonminus + \Pproton \\
      \APnulepton + \Pproton \rightarrow \Pleptonplus + \Pneutron
    \end{align*}


\end{document}